\documentclass[11pt,a4paper]{article}
\usepackage{cafestyle}
\setlength{\parindent}{0pt}
\setlength{\parskip}{0.8em plus 0.2em minus 0.1em}

% \analysis{number}{label}{question}{null}{alternative}{methods}
\newcommand{\analysis}[6]{%
  \par\smallskip\noindent\textbf{A#1\enspace #2.}\enspace #3\enspace
  \Ho: #4\enspace \Hone: #5\enspace \textit{Methods:} #6\par}
\title{\textbf{Semester Project: Data Analysis}\\[0.4ex]\large The Prague caf\'e data: overview, research questions, tests and reporting standards}
\author{Data Analytics}
\date{\today}

\begin{document}
\maketitle

\begin{abstract}
\noindent This document describes the data analysis part of the semester project. You analyse data on caf\'es in Prague: the prices of four coffees, the Google rating and number of ratings, seating, the distance from the Orloj, the district and the caf\'e type. You work in a batch of two students. Each batch answers its own questions with the same ten methods. You submit three R scripts and a \LaTeX\ report, and you give a presentation of 10 to 15 minutes.
\end{abstract}

\clearpage
\tableofcontents

\clearpage

\section*{Project Overview}
The semester project is worth 50 percentage points of your grade and has three stages:

\begin{itemize}
	\item \textbf{Data Collection} (20~pp). You collect real-world data on coffee prices in Prague. See \hyperref[sec:collection]{\textcolor{blue}{Section~\ref*{sec:collection}}}.
	\item \textbf{Data Analysis} (15~pp). You clean, describe and analyse the collected data, and report your findings in a written report and a presentation. See \hyperref[sec:analysis]{\textcolor{blue}{Section~\ref*{sec:analysis}}}.
	\item \textbf{Analysis Defense} (15~pp). You show that you understand the analysis you submitted. See \hyperref[sec:defense]{\textcolor{blue}{Section~\ref*{sec:defense}}}.
\end{itemize}

The project requires R and RStudio, GitHub and \LaTeX\ (for both typesetting and presentation), so you will need to become familiar with them along the way.

The syllabus expects 2 hours of work outside class for every class hour, which makes 84 hours per semester. As the project is 50\% of your grade, plan for about 40 hours of work, roughly one full working week.

\clearpage
\section{Data Collection}\label{sec:collection}
The semestral project starts with a data collection you are going to do. 

\subsection*{General rules}

\begin{itemize}
\item Each of you must collect 20 unique places -- coordinate between yourself the locations so you don't end up with the same spots, duplicates do not count. If you collect more than 20, good (but no extra credit).
\item Prefer cafes, not restaurants that make coffee.
\item Only 2 of the 20 can be an international chain such as Starbucks, Costa, etc.
\item Introduce yourself, inform staff what you are doing (see below).
\end{itemize}

\subsection*{Introduce yourself!}

\begin{itemize}
\item Before collecting data at each caf\'e, briefly introduce yourself to a staff member and explain that you're a university student collecting publicly available information (prices, seating, and similar details) for a data analytics course assignment -- not a mystery shopper, journalist, or inspector, and not there to critique the business.
\item Ask before photographing the menu or price board, and never photograph other customers or staff.
\item Keep the interaction short and don't take up staff time during a busy rush; if a caf\'e is clearly slammed, come back later or move on to another location rather than pushing.
\item You are not obligated to make a purchase to collect this data.
\item If staff would rather you not collect data there, thank them and leave -- respect the decision and pick a different caf\'e rather than arguing or persisting.
\item Throughout, be polite, honest about what you're doing and why, and remember that how you behave reflects on both you and the university, since these are real small businesses being asked for a small favor, not test subjects.
\end{itemize}

You can use, for instance, this:
\begin{quote}
\itshape ``Hi, I'm a university student working on a class project for my Data Analytics course. We're collecting some publicly available information about caf\'es around Prague -- things like drink prices and seating -- it'll just take a minute, and I wanted to let you know before I look around and note a few things. Is that okay?''
\end{quote}

\subsection*{How to fill the sheet}

\begin{description}
\item[Your Name] Full name as enrolled, so rows can be matched to the right student.
\item[Date and time of visit] Date and local time you were actually there, in a consistent format (M/D/YYYY H:MM). Matters for spotting outliers and confirming coverage.
\item[District (number)] The official Prague district (e.g.\ ``Prague 2''), not the postal code.
\item[District (name)] The commonly used neighborhood name (e.g.\ ``Vinohrady''), not a restatement of the number.
\item[Cafe name] As shown on the caf\'e's signage or its Google Maps listing.
\item[Cafe address (Google Maps)] Full address exactly as Google Maps displays it.
\item[Coordinates] Find the place on Google Maps. Right-click the pin, then click the coordinates that appear at the top of the menu (auto-copies). Paste as-is, one cell, comma-separated.
\item[Link to GMaps] Share $\rightarrow$ ``Copy link,'' not ``Embed a map.''
\item[Espresso / Cappuccino / Flat White / Filter-Batch Price in CZK (dine-in)] Dine-in price only, as listed on the menu or POS display. If the caf\'e doesn't serve that drink, write NA -- never 0 or a blank cell, since those would look like ``free'' or ``forgot to check.''
\item[Payment] Exactly one of: Cash only / Card only / Cash \& Card (use the dropdown, don't free-type).
\item[Seating Places Inside / Outside] Count individual seats, not tables. Write 0 if there genuinely are none -- don't leave blank. Ideally, ask the staff.
\item[Rating (Google) / Number of ratings (Google)] Copied directly from the caf\'e's Maps listing at the time of your visit.
\item[Cafe Type] Exactly one of:
  \begin{itemize}
  \item Independent: single, standalone location; no other branches under that name anywhere.
  \item Local chain: 2+ locations, all within Czech Republic.
  \item International chain: 2+ locations, present in more than one country.
  \end{itemize}
  Ideally, ask the staff. Alternatively, check via Google Maps or the caf\'e's website/Instagram: if multiple pins/locations show up under the same brand, it's a chain. Use the dropdown, don't free-type.
\item[Is Roastery (0/1)] 1 if it roasts its own beans on-site or is branded as a roastery, 0 otherwise. Always fill this in --- 0 is a real answer, not ``unknown.''
\item[Photo Folder] Link to (or your named subfolder in) the shared class Drive where you've uploaded a photo of the menu/price board, as backup evidence for your prices.
\item[Note] Free text for anything that affects how your data should be read (e.g.\ ``no regular milk, oat milk only'' -- explains why a price might not be directly comparable).
\end{description}

\clearpage
\section{Data Analysis}\label{sec:analysis}

In this part of the project, you analyse the data. You will:

\begin{itemize}
	\item clean the data (see \hyperref[sec:rules]{\textcolor{blue}{subsection~\ref*{sec:rules}}}),
	\item analyse the data according to the research questions and hypotheses assigned to you (see \hyperref[sec:batches]{\textcolor{blue}{subsection~\ref*{sec:batches}}}),
	\item write a report on your findings (see \hyperref[sec:deliverables]{\textcolor{blue}{subsection~\ref*{sec:deliverables}}}),
	\item present your findings (see \hyperref[sec:deliverables]{\textcolor{blue}{subsection~\ref*{sec:deliverables}}}).
\end{itemize}

You do this in a team of one or two classmates of your choice, according to the lecturer's instructions. In the team, you may cooperate on everything and use AI as much as you want. Each of you, however, submits your own complete version of the analysis in your own GitHub repository. The only exception is the presentation: it deals only with selected topics, so each of you presents half of the findings if you are two, a third if you are three, and you decide the division among yourselves.

Your understanding of the code and of the whole analysis will be checked during the defense (see Section~\ref{sec:defense}), where I may ask you about anything in your analysis, including the parts that you did not present. So treat cooperation and AI as a learning process that helps you handle and understand everything, not as something that does the work for you without your understanding it.

\subsection{Data and general rules}\label{sec:rules}

\subsubsection*{The variables}

Each row of the data file is one caf\'e. The file contains the date and time of the visit, the district (number and name), the name and address of the caf\'e, its coordinates and a Google Maps link. Prices are in Czech crowns for a dine-in drink with normal milk: espresso, cappuccino, flat white and filter or batch coffee (NA if not offered). Payment methods are given as text. Seats inside and outside are counts. The Google rating is on a scale from 1 to 5, and the number of ratings is a count that is sometimes written with a thousands separator. The caf\'e type is independent, local chain or international chain. The roastery indicator is 1 for a roastery and 0 otherwise, and some cells are blank. The distance from the Orloj (the astronomical clock in the centre of Prague) is in kilometres and is calculated from the latitude and longitude, which are also in the file. The photo folder and the free-text note are not used in the analyses, but the notes sometimes explain unusual prices. The data were entered by hand, so labels and formats are not uniform and some cells contain more than one value.

\subsubsection*{Cleaning and general rules}

Clean the raw data with your own script. The script reads the raw file, which you do not edit by hand, performs every cleaning step in code and writes the cleaned data set that you use in the analysis. At a minimum it drops empty rows, harmonises the district labels (for example ``prague 1'' and ``prauge 1''), removes the thousands separator from the number of ratings and sets a blank caf\'e type to missing. Cells with several prices (for example ``85, 95, 105'' or ``65/40'') need an explicit rule. By default use the first listed price, which is the dine-in or smallest size. Describe each cleaning decision in the report, together with the number of caf\'es it affects, because different decisions can change the results.

The number of ratings and the number of seats are strongly right-skewed, and one caf\'e has no ratings. In correlations and regressions use both as $\log(x+1)$, whether as outcome or as predictor. Where a question compares raw values, for example seats inside against seats outside, use the raw counts and handle the skewness through your choice of method. District comparisons use only districts with at least 12 caf\'es (Prague 1, 2, 3, 5, 7 and 8). If a question gives no benchmark for the one-sample test, choose one before you look at the data and state in the report where it came from. Write down your hypotheses, including the direction of any one-tailed alternative, before you run any test.

Draw the figures before you run the corresponding test, and check residual plots for every regression. Label results that you did not plan as exploratory. Your grade does not depend on statistical significance. A clear null result with a good interpretation is a full answer.

\subsection{Reporting standards}

\subsubsection*{General conventions}

Use $\alpha=.05$ and state it once in the methods section. Report test statistics with degrees of freedom, as in $t(df)$, $F(df_1,df_2)$ or $\chi^2(df,N)$, and give Welch degrees of freedom with one decimal. Give exact $p$-values to three decimals without a leading zero (\pr\,=\,.034), and write very small values as \plt{.001}. Round estimates and test statistics to two decimals and prices to one decimal, and give a confidence interval the same number of decimals as its estimate. Set statistical symbols in italics and give every estimate its unit (CZK, km, rating points). Report every number with the sample size it is based on, and state for each analysis how many caf\'es were lost to missing values. The numbers in the examples below are invented.

Report a result as an estimate, its uncertainty and its size, in this order, and only then as a test. Say what the estimate means in the units of the data and whether a difference of that size is practically relevant.

\subsubsection*{Describing the data}

Describe the data before you run any test, in the second script. Summarise the description in the data section of the report and give it in full in an appendix. For every variable you use, report the number of observations and of missing values. For numeric variables report the mean, standard deviation, median, interquartile range and range, and for categorical variables the counts and percentages. Plot every numeric variable with a histogram or density plot and a boxplot, and plot the groups to be compared side by side. Check normality within each group (for a regression, of the residuals) with a normal Q--Q plot and the skewness. You may add the Shapiro--Wilk test, but the plots take priority, because the test has little power in small samples and flags unimportant deviations in large ones. State the result of the check in the report and which of the two methods named for a question you therefore treat as the primary one. Report the other as a check.

\subsubsection*{Means and medians}

Report every mean with its 95\% confidence interval and its sample size, for example \meanci{82.4}{79.8}{85.0}{CZK}, $n=312$. Give the standard deviation in the descriptive table. For skewed variables report the median and interquartile range next to the mean. Where the plots suggest that the normal approximation is doubtful, use a bootstrap confidence interval (percentile or BCa, at least 5{,}000 resamples, seed stated) instead of the $t$-based interval. Show confidence intervals in the figures as well as in the text.

\subsubsection*{Differences between means}

State the direction of the subtraction in words (``Prague 1 minus the rest''). Report the difference with its 95\% confidence interval and a standardised effect size with its own confidence interval: Cohen's $d$ for two independent groups (Hedges' $g$ if a group has fewer than 20 observations) and $d_z$ for paired data, which is the mean difference divided by the standard deviation of the differences. Example: \diffci{4.1}{1.2}{7.0}{CZK}{0.31}{0.09}{0.53}. You may mention Cohen's benchmarks (0.2, 0.5, 0.8) once, but they do not replace an interpretation in crowns. For a one-tailed alternative, still show the two-sided interval and state that the $p$-value is one-tailed and was fixed in advance. If you run a rank test as well as a $t$-test, report the Hodges--Lehmann shift or the rank-biserial correlation for it, with a confidence interval.

\subsubsection*{Three or more groups}

Report Welch's $F(df_1,df_2)$ with $\eta^2$ (or $\omega^2$) and its confidence interval, or the Kruskal--Wallis $H(df)$ with $\varepsilon^2$. Follow a significant result with pairwise comparisons (Tukey, Games--Howell, or Dunn with Holm adjustment) and give for each pair the mean difference, the adjusted $p$-value and the confidence interval. Show the group means with their confidence intervals in one figure.

\subsubsection*{Two categorical variables}

Report $\chi^2(df,N)$, the counts and row percentages, the smallest expected count, and Cram\'er's $V$ with its confidence interval. If any expected count is below 5, use Fisher's exact test and, for $2\times2$ tables, report the odds ratio with its confidence interval. Standardised residuals above 2 in absolute value show which cells drive the association.

\subsubsection*{Correlation}

Show the scatterplot first. Report the coefficient with its 95\% confidence interval and $n$, for example $r=-.21$, \CI{-.30}{-.12}, and the squared coefficient as the share of shared variance. Report Spearman's $\rho$ and Kendall's $\tau$ next to Pearson's $r$, with a bootstrap confidence interval for at least one of them. A correlation is not an effect, so do not suggest a cause.

\subsubsection*{Regression}

Export the regression table from R. It contains the number of observations, each coefficient with its robust standard error and 95\% confidence interval, $R^2$, adjusted $R^2$ and the residual standard error. State the slope in units of the outcome per unit of the predictor, and for a log outcome as a percentage change (a coefficient $b$ means a change of $100\,(e^{b}-1)\%$). Inspect the residual-versus-fitted plot, the normal Q--Q plot and the leverage plot, name influential observations and say whether your conclusion holds without them. Report a nested model comparison as $F(df_1,df_2)$ and the change in $R^2$, and for the multiple regression the variance inflation factors. State how much the slope of interest changed when the second predictor entered, in absolute and relative terms.

\subsubsection*{Interactions}

Report the interaction coefficient with its 95\% confidence interval, the nested $F$-test against the model without it and the change in $R^2$. Then state the slope of interest for each group with its own confidence interval, because the interaction coefficient is only the difference between two slopes. Show the fitted lines by group over the data. A non-significant interaction does not show that the slopes are equal, so say how large a difference the interval still allows.

\subsubsection*{Logistic regression}

The main result is the average marginal effect of the predictor of interest, stated as a change in probability in percentage points with its 95\% confidence interval. For example: ``each additional kilometre from the Orloj lowers the probability of a cappuccino above the median price by 5.1 percentage points, \CI{0.2}{10.0}, other things equal.'' You may add the odds ratio with its confidence interval as a check. Report the number of observations, the number of events and McFadden's pseudo-$R^2$, and show a predicted-probability plot with a confidence band. If you create the outcome by a median split, state the median.

\subsubsection*{Many tests and exploratory results}

Count the tests in your report. Adjust tests that answer one family of questions (for example all pairwise district comparisons) with Holm's method. Label anything that you did not plan in advance as exploratory. A single significant result among many is a hypothesis for further work, not a finding.

\subsubsection*{Figures, tables and interpretation}

Every figure shows the data (points, distribution or both) with the estimate and its confidence interval, has labelled axes with units and has a caption that can be understood without the text. Every table has a caption too. Produce tables with the R script and include them with \verb|\input|. Do not type them by hand, and use \texttt{booktabs} rules. The discussion covers three limits: the sample is a convenience sample and not a random sample of Prague caf\'es, the data are observational, and prices of caf\'es in the same district are not independent. Write ``in this sample'' unless you can defend a wider claim.

\clearpage
\subsection{The five batches}\label{sec:batches}

\subsubsection{Batch 1: cappuccino price and location}

The outcome is the cappuccino price in CZK. The question is whether location (distance from the Orloj, Prague 1, district) and caf\'e type explain it.

\analysis{1}{One sample}{Is the average Prague cappuccino price different from your own prior guess?}{the mean equals the benchmark you chose in advance.}{the mean differs from it.}{one-sample $t$-test, Wilcoxon signed-rank test, bootstrap confidence interval.}
\analysis{2}{Two groups}{Is a cappuccino dearer in Prague 1 than in the rest of Prague?}{mean prices are equal.}{Prague 1 is higher (one-tailed).}{Welch $t$-test, Wilcoxon rank-sum test, Cohen's $d$.}
\analysis{3}{Paired}{How much more does a cappuccino cost than an espresso in the same caf\'e?}{the mean difference is zero.}{the cappuccino is dearer (one-tailed).}{paired $t$-test, Wilcoxon signed-rank test.}
\analysis{4}{Several groups}{Do cappuccino prices differ across districts 1, 2, 3, 5, 7 and 8?}{all district means are equal.}{at least one differs.}{Welch ANOVA, Kruskal--Wallis test, Games--Howell or Dunn post hoc test.}
\analysis{5}{Two categorical variables}{Is the caf\'e type (independent, local chain, international chain) related to being in Prague 1?}{type and location are independent.}{they are associated.}{chi-square test or Fisher's exact test, Cram\'er's $V$, standardised residuals. The international chain group has only eleven caf\'es, so check the expected counts before choosing.}
\analysis{6}{Correlation}{Does the price fall with the distance from the Orloj?}{$\rho=0$.}{$\rho<0$, because the centre is dearer (one-tailed).}{Pearson, Spearman and Kendall coefficients, bootstrap confidence interval.}
\analysis{7}{Simple regression}{By what percentage does the price of a cappuccino change with each kilometre from the Orloj, and is the relationship better described in crowns per kilometre?}{the slope is zero.}{the slope is not zero.}{simple regression of $\log(\text{price})$ on distance (log-linear), compared with price on distance in crowns, residual plots, robust standard errors. With the outcome in logs, $100\,(e^{b}-1)$ is the percentage change per kilometre, and the statement needs a sentence on what the data can and cannot say about why.}
\analysis{8}{Multiple regression}{Does the distance slope survive once caf\'e type is in the model?}{caf\'e type adds nothing.}{caf\'e type improves the model.}{multiple regression of price on distance and type, nested $F$-test, variance inflation factors.}
\analysis{9}{Interaction}{Does the price fall with the distance from the Orloj in the same way for independent caf\'es and for chains (local and international together)?}{the interaction coefficient is zero, so the distance slope is the same in both groups.}{the interaction coefficient is not zero.}{multiple regression of price on distance, chain indicator and their interaction, nested $F$-test against the model without the interaction, the distance slope in each group with its confidence interval, plot of the fitted lines.}
\analysis{10}{Logistic regression}{By how many percentage points does the probability that a cappuccino costs more than the median price (83 CZK) change with each kilometre from the Orloj?}{the distance coefficient is zero.}{the distance coefficient is not zero.}{logit of the indicator on distance and caf\'e type, average marginal effect of distance in percentage points, predicted-probability plot.}

\clearpage
\subsubsection{Batch 2: espresso price and caf\'e type}

The outcome is the espresso price in CZK. The question is what kind of caf\'e charges more.

\analysis{1}{One sample}{Is the average espresso price different from your own prior guess?}{the mean equals the benchmark you chose in advance.}{the mean differs from it.}{one-sample $t$-test, Wilcoxon signed-rank test, bootstrap confidence interval.}
\analysis{2}{Two groups}{Do chains (local and international together) charge more for espresso than independents?}{mean prices are equal.}{chains are higher (one-tailed).}{Welch $t$-test, Wilcoxon rank-sum test, Cohen's $d$.}
\analysis{3}{Paired}{How much more does a flat white cost than an espresso in the same caf\'e?}{the mean difference is zero.}{the flat white is dearer (one-tailed).}{paired $t$-test, Wilcoxon signed-rank test.}
\analysis{4}{Several groups}{Does the espresso price differ across the three caf\'e types?}{all type means are equal.}{at least one differs.}{Welch ANOVA, Kruskal--Wallis test, Games--Howell or Dunn post hoc test. The international chain group has only eleven caf\'es and a much wider spread, so Welch and Games--Howell are safer than the classic ANOVA and Tukey.}
\analysis{5}{Two categorical variables}{Is being a roastery related to caf\'e type?}{the two are independent.}{they are associated.}{Fisher's exact test because the cells are small, Cram\'er's $V$.}
\analysis{6}{Correlation}{Is the espresso price related to the Google rating?}{$\rho=0$.}{$\rho\neq0$ (two-tailed, no strong prior).}{Pearson, Spearman and Kendall coefficients, bootstrap confidence interval.}
\analysis{7}{Simple regression}{How many crowns does the espresso change per rating point?}{the slope is zero.}{the slope is not zero.}{simple regression of espresso price on rating, residual plots, robust standard errors.}
\analysis{8}{Multiple regression}{Does the rating slope change once caf\'e type is in the model?}{caf\'e type adds nothing.}{caf\'e type improves the model.}{multiple regression of price on rating and type, nested $F$-test, variance inflation factors.}
\analysis{9}{Interaction}{Is the relationship between the espresso price and the rating the same for independent caf\'es and for chains (local and international together)?}{the interaction coefficient is zero, so the rating slope is the same in both groups.}{the interaction coefficient is not zero.}{multiple regression of espresso price on rating, chain indicator and their interaction, nested $F$-test against the model without the interaction, the rating slope in each group with its confidence interval, plot of the fitted lines.}
\analysis{10}{Logistic regression}{By how many percentage points does the probability that an espresso costs more than the median price (65 CZK) change with each kilometre from the Orloj?}{the distance coefficient is zero.}{the distance coefficient is not zero.}{logit of the indicator on distance and caf\'e type, average marginal effect of distance in percentage points, predicted-probability plot.}

\clearpage
\subsubsection{Batch 3: flat white price and the venue}

The outcome is the flat white price in CZK. The question is whether the size and setting of the caf\'e (seats, outdoor seating, distance) go with price.

\analysis{1}{One sample}{Is the average flat white price different from your own prior guess?}{the mean equals the benchmark you chose in advance.}{the mean differs from it.}{one-sample $t$-test, Wilcoxon signed-rank test, bootstrap confidence interval.}
\analysis{2}{Two groups}{Do caf\'es with outdoor seating charge more for a flat white?}{mean prices are equal.}{caf\'es with outdoor seating are higher (one-tailed).}{Welch $t$-test, Wilcoxon rank-sum test, Cohen's $d$.}
\analysis{3}{Paired}{How much more does a flat white cost than a cappuccino in the same caf\'e?}{the mean difference is zero.}{the flat white is dearer (one-tailed).}{paired $t$-test, Wilcoxon signed-rank test.}
\analysis{4}{Several groups}{Does the flat white price differ across the three caf\'e types?}{all type means are equal.}{at least one differs.}{Welch ANOVA, Kruskal--Wallis test, Games--Howell or Dunn post hoc test.}
\analysis{5}{Two categorical variables}{Is outdoor seating related to being in Prague 1?}{the two are independent.}{they are associated.}{chi-square test, Cram\'er's $V$, standardised residuals.}
\analysis{6}{Correlation}{Is the flat white price related to the number of seats inside?}{$\rho=0$.}{$\rho>0$, because bigger caf\'es charge more (one-tailed).}{Pearson on $\log(\text{seats}+1)$, Spearman and Kendall coefficients, bootstrap confidence interval. One caf\'e has 260 seats, so work on the log scale and show that the result does not depend on that row.}
\analysis{7}{Simple regression}{How many crowns does a flat white change when the number of seats doubles?}{the slope is zero.}{the slope is not zero.}{simple regression of price on $\log(\text{seats}+1)$, residual plots, robust standard errors.}
\analysis{8}{Multiple regression}{Does the seats slope change once distance is in the model?}{distance adds nothing.}{distance improves the model.}{multiple regression of price on log seats and distance, nested $F$-test, variance inflation factors.}
\analysis{9}{Interaction}{Does the effect of caf\'e size on the flat white price differ between Prague 1 and the rest of Prague?}{the interaction coefficient is zero, so the seats slope is the same in both areas.}{the interaction coefficient is not zero.}{multiple regression of price on $\log(\text{seats}+1)$, a Prague 1 indicator and their interaction, nested $F$-test against the model without the interaction, the seats slope in each area with its confidence interval, plot of the fitted lines.}
\analysis{10}{Logistic regression}{By how many percentage points does the probability that a caf\'e has outdoor seating change with each kilometre from the Orloj?}{the distance coefficient is zero.}{the distance coefficient is not zero.}{logit of an indicator for at least one outdoor seat on distance and $\log(\text{seats}+1)$, average marginal effect of distance in percentage points, predicted-probability plot.}

\clearpage
\subsubsection{Batch 4: Google rating}

The outcome is the Google rating (3.6 to 5.0). The question is what goes with a better-rated caf\'e. Ratings are bounded and skewed, so the rank-based tests carry more weight in this batch.

\analysis{1}{One sample}{Is the average rating different from 4.5, the level one would call a good caf\'e?}{the mean equals 4.5.}{the mean differs from 4.5.}{one-sample $t$-test, Wilcoxon signed-rank test, bootstrap confidence interval.}
\analysis{2}{Two groups}{Do roasteries have a different rating from other caf\'es?}{mean ratings are equal.}{they differ (two-tailed).}{Welch $t$-test, Wilcoxon rank-sum test, Cohen's $d$.}
\analysis{3}{Paired}{Do caf\'es have more seats inside than outside?}{the mean difference is zero.}{there are more seats inside (one-tailed).}{paired $t$-test on the differences, Wilcoxon signed-rank test. Seat counts are very skewed, so check the paired differences with a plot and compare the $t$-test with the Wilcoxon result.}
\analysis{4}{Several groups}{Does the rating differ across the three caf\'e types?}{all type distributions are equal.}{at least one differs.}{Kruskal--Wallis test, Dunn post hoc test with Holm adjustment, Welch ANOVA as a check.}
\analysis{5}{Two categorical variables}{Is a high rating (at or above the median, 4.6) related to caf\'e type?}{the two are independent.}{they are associated.}{Fisher's exact test because international chains are few, Cram\'er's $V$.}
\analysis{6}{Correlation}{Is the rating related to how many people rated the caf\'e?}{$\rho=0$.}{$\rho\neq0$.}{Pearson on $\log(\text{ratings}+1)$, Spearman and Kendall coefficients, bootstrap confidence interval.}
\analysis{7}{Simple regression}{How many rating points does the rating change per doubling of the number of ratings?}{the slope is zero.}{the slope is not zero.}{simple regression of rating on $\log(\text{ratings}+1)$, residual plots, robust standard errors.}
\analysis{8}{Multiple regression}{Does the slope change once the cappuccino price is in the model?}{price adds nothing.}{price improves the model.}{multiple regression of rating on log ratings and cappuccino price, nested $F$-test, variance inflation factors.}
\analysis{9}{Interaction}{Is the relationship between the rating and the number of ratings the same in Prague 1 and in the rest of Prague?}{the interaction coefficient is zero, so the slope is the same in both areas.}{the interaction coefficient is not zero.}{multiple regression of rating on $\log(\text{ratings}+1)$, a Prague 1 indicator and their interaction, nested $F$-test against the model without the interaction, the slope in each area with its confidence interval, plot of the fitted lines.}
\analysis{10}{Logistic regression}{By how many percentage points does the probability of a rating at or above the median (4.6) change with each kilometre from the Orloj?}{the distance coefficient is zero.}{the distance coefficient is not zero.}{logit of the indicator on distance and $\log(\text{ratings}+1)$, average marginal effect of distance in percentage points, predicted-probability plot.}

\clearpage
\subsubsection{Batch 5: popularity and size}

The outcome is popularity, measured by the number of Google ratings on the log scale. The question is where the busy caf\'es are.

\analysis{1}{One sample}{Is the typical number of seats inside different from 30?}{the mean equals 30.}{the mean differs from 30.}{one-sample $t$-test, Wilcoxon signed-rank test, bootstrap confidence interval; the data are skewed, so lead with the bootstrap.}
\analysis{2}{Two groups}{Do caf\'es in Prague 1 have more ratings than caf\'es elsewhere?}{mean log ratings are equal.}{Prague 1 is higher (one-tailed).}{Welch $t$-test, Wilcoxon rank-sum test, Cohen's $d$.}
\analysis{3}{Paired}{Does a filter coffee cost the same as a cappuccino in the same caf\'e?}{the mean difference is zero.}{the prices differ (two-tailed).}{paired $t$-test, Wilcoxon signed-rank test. Only about 54 caf\'es list a filter price.}
\analysis{4}{Several groups}{Does the number of ratings differ across districts 1, 2, 3, 5, 7 and 8?}{all district means of log ratings are equal.}{at least one differs.}{Welch ANOVA, Kruskal--Wallis test, Games--Howell or Dunn post hoc test.}
\analysis{5}{Two categorical variables}{Is a high cappuccino price (above the median, 83 CZK) related to being in Prague 1?}{the two are independent.}{they are associated.}{chi-square test, Cram\'er's $V$, standardised residuals.}
\analysis{6}{Correlation}{Does popularity fall with the distance from the Orloj?}{$\rho=0$.}{$\rho<0$ (one-tailed).}{Pearson on $\log(\text{ratings}+1)$, Spearman and Kendall coefficients, bootstrap confidence interval.}
\analysis{7}{Simple regression}{By what percentage does the number of ratings change per kilometre?}{the slope is zero.}{the slope is not zero.}{simple regression of log ratings on distance (log-linear), residual plots, robust standard errors. With the outcome in logs, $e^{b}-1$ gives the percentage change per kilometre, and the statement needs a sentence on what the data can and cannot say about why.}
\analysis{8}{Multiple regression}{Does the distance slope change once caf\'e type is in the model?}{caf\'e type adds nothing.}{caf\'e type improves the model.}{multiple regression of log ratings on distance and type, nested $F$-test, variance inflation factors.}
\analysis{9}{Interaction}{Does popularity fall with the distance from the Orloj in the same way for caf\'es with and without outdoor seating?}{the interaction coefficient is zero, so the distance slope is the same in both groups.}{the interaction coefficient is not zero.}{multiple regression of $\log(\text{ratings}+1)$ on distance, an outdoor-seating indicator and their interaction, nested $F$-test against the model without the interaction, the distance slope in each group with its confidence interval, plot of the fitted lines.}
\analysis{10}{Logistic regression}{By how many percentage points does the probability of having more than the median number of Google ratings (454) change with each kilometre from the Orloj?}{the distance coefficient is zero.}{the distance coefficient is not zero.}{logit of the indicator on distance and caf\'e type, average marginal effect of distance in percentage points, predicted-probability plot.}

\subsection{Deliverables and grading}\label{sec:deliverables}

This part of the semester project is worth 15 points, divided equally among three components. Your R code (5 points) consists of three scripts that you run in this order: \texttt{01\_cleaning.R} turns the raw data file into the data set for the analysis, \texttt{02\_descriptives.R} describes the data and checks distributions and normality, and \texttt{03\_tests.R} runs the tests and regressions. Together the scripts reproduce every number, table and figure in your report and slides. Your \LaTeX\ report (5 points) is written in your own words from the supplied template and contains tables and figures exported from R, as described below. Your presentation (5 points) has at least five slides, lasts between 10 and 15 minutes and is spoken, not read.

\subsubsection*{Tables and figures}

Produce every table and figure of the report in R. Do not type or draw anything by hand. For each of the ten analyses, export at least one table and one figure, named by analysis (for example \texttt{A2\_table.tex} and \texttt{A2\_figure.pdf}), to \texttt{report/tables} and \texttt{report/figures}. The table gives the estimates with their confidence intervals, the effect size, the test statistic with its degrees of freedom and the $p$-value. The figure shows the data with the estimate and its confidence interval. Every table and figure has a caption, which you add in the \LaTeX\ source around the exported file and not in the R output. These are minimums: the data description, the diagnostics and the robustness checks add more. The main text of the report contains the key figure of each analysis and one summary table of the main results of all ten analyses. The other files go in the appendices.

\subsubsection*{Report}

The report has no page limit. It contains everything a reader needs to understand what you did: the data and cleaning, the methods, the results and their interpretation, and the limits of the analysis. The reader should not need your code. The main text contains only the main results: for each analysis the estimate with its confidence interval, the effect size, the supporting test, one key figure and the interpretation. Everything that checks or supports these results goes in the appendices and is referred to from the main text: the full descriptive statistics, the distribution and normality checks, the regression diagnostics, the alternative tests and robustness analyses, the full post hoc tables and the reproducibility information.

\subsubsection*{Grading criteria}

Your code is graded on whether it runs without errors from the raw file through the three scripts in order on a fresh copy of the repository, whether all cleaning and analysis steps are done in code, whether the scripts are organised in commented blocks with one block per analysis, whether paths are relative and random procedures use a stated seed, and whether the exported tables and figures are the ones used in the report. Your report is graded on whether a reader can understand what you did without the code, whether hypotheses, estimates, confidence intervals, effect sizes and test statistics follow the reporting standards, whether the interpretation is correct and given in the units of the data, whether main text and appendices are separated as described, whether every table and figure is captioned and referred to, and whether the limits are stated. Your presentation is graded on whether it concentrates on the findings and their interpretation, names the test that supports each finding, is spoken and not read, has at least five slides and keeps to 10 to 15 minutes.

\subsubsection*{Distribution of points}

The 5 points for the code are distributed as follows: 1.5 points for code that runs without errors from the raw file through the three scripts in order, 1 point for complete cleaning done in code with every decision recorded, 1.5 points for the ten analyses using the specified methods, each in its own commented block, 0.5 points for tables and figures exported by script and matching the report, and 0.5 points for relative paths, stated seeds and a readable structure. The 5 points for the report are distributed as follows: 1 point for a clear account of the data, cleaning and methods, 1 point for numbers reported in the required format, 1.5 points for a correct interpretation in the units of the data, 1 point for the separation of main text and appendices and for captioned tables and figures, and 0.5 points for the limits of the analysis and the labelling of exploratory results. The 5 points for the presentation are distributed as follows: 2 points for clear and correct findings and interpretations, 1 point for naming the supporting test of each finding together with a good plot or estimate, 1 point for delivery (the talk is spoken and not read, and the findings are shared fairly between you and your partner), 0.5 points for at least five slides and a length of 10 to 15 minutes, and 0.5 points for the quality of the slides. These points reward completeness and sound reporting.

\subsubsection*{Repository}

Create your own GitHub repository containing your own code and report, and send the link to the lecturer. You do not share a repository with your partner. The repository has three folders. The folder \texttt{analysis} holds the raw data, the three scripts and the cleaned data set. The folder \texttt{report} holds the \LaTeX\ source, the PDF and the exported tables and figures. The folder \texttt{presentation} holds the slides as a PDF and in their source format. Use paths relative to the repository, so that the analysis can be reproduced after cloning it. Apart from the raw data, do not upload any file that the scripts cannot regenerate.

\subsubsection*{Presentation}

Present the findings, not the procedure. Each of you presents about half of the findings of your batch, and you decide how to divide them. For every finding show the key plot or estimate with its confidence interval and effect size, give the interpretation in plain words and name the test that supports it. Do not present the details of the tests (assumptions, alternatives, computation of the statistic).

\clearpage
\section{Analysis Defense}\label{sec:defense}

The analysis defense is worth 15~pp. It is an oral examination in which you show that you understand the analysis you submitted. It is based on your own code, your own report and your own presentation, and each of you is examined individually, also when you worked on the same questions with your partner.

You should be able to explain every line of your three scripts and every decision in them: how you cleaned the data and what each cleaning step changes, why you chose each test, which assumptions it makes and how you checked them. You should also be able to explain every number in your report and your presentation, including the confidence intervals, the effect sizes and the $p$-values, and to interpret them in the units of the data without claiming more than the data support.

\end{document}
